\documentclass[12pt,a4paper]{article}

\usepackage[T1]{fontenc}
\usepackage[utf8]{inputenc}
\usepackage[brazil]{babel}
\usepackage{lmodern}
\usepackage{geometry}
\usepackage{hyperref}
\usepackage{enumitem}
\usepackage{array}

\geometry{margin=2.5cm}
\hypersetup{
  colorlinks=true,
  linkcolor=blue,
  urlcolor=blue
}

\title{Plano de Aula: Aula 5 -- Prototipagem de Interfaces de Software}
\date{}

\begin{document}
\maketitle

\noindent
\textbf{Disciplina:} Engenharia de Software (Curso Técnico em Desenvolvimento de Sistemas)\\
\textbf{Duração:} 2 horas (120 minutos)\\
\textbf{Modalidade:} Presencial\\
\textbf{Materiais da aula:} \href{aula-05.html}{slides (HTML)} \quad \href{aula-05.pdf}{ementa (PDF)}\\
\textbf{Livro-base (referência):} \href{../99-Materiais/Engenharia%20de%20software%20projetos%20e%20processos%20by%20Wilson%20de%20P%C3%A1dua%20Paula%20Filho%20.epub.pdf}{Engenharia de software: projetos e processos (PDF)}

\section{Competências e Habilidades}
\textbf{Competências}
\begin{itemize}[leftmargin=*,itemsep=0.25em]
  \item Aplicar princípios da engenharia de software para reduzir retrabalho por validação tardia.
  \item Planejar e representar interfaces e fluxos antes da implementação.
\end{itemize}

\textbf{Habilidades}
\begin{itemize}[leftmargin=*,itemsep=0.25em]
  \item Construir fluxo de telas e wireframes.
  \item Diferenciar protótipo de baixa, média e alta fidelidade.
  \item Definir perguntas de validação e critérios de sucesso.
\end{itemize}

\section{Objetivos da Aula}
\begin{itemize}[leftmargin=*,itemsep=0.25em]
  \item Explicar por que prototipar reduz custo de mudança e risco.
  \item Criar fluxo de navegação mínimo para uma funcionalidade.
  \item Produzir wireframe da tela principal de uma solução.
  \item Preparar uma validação rápida com perguntas e observação.
\end{itemize}

\section{Assuntos Abordados no Dia}
\begin{itemize}[leftmargin=*,itemsep=0.25em]
  \item Protótipo como ferramenta de engenharia: reduzir incerteza e retrabalho.
  \item Fidelidade: baixa, média e alta; quando usar cada uma.
  \item Fluxo de telas e estados (sucesso, erro, vazio).
  \item Wireframe: hierarquia, componentes, mensagens e validações.
  \item Validação: perguntas, observação e registro de feedback.
\end{itemize}

\section{Metodologia}
Aula com demonstrações rápidas no quadro (exemplo de fluxo e wireframe). Atividade em grupo para desenhar protótipo em papel (baixa fidelidade) e preparar validação. O professor aplica micro-feedback por ciclos: rascunho $\rightarrow$ perguntas $\rightarrow$ refinamento.

\section{Materiais e Recursos}
\begin{itemize}[leftmargin=*,itemsep=0.25em]
  \item Quadro e marcadores.
  \item Papel A4, lápis, canetas, post-its (opcional).
  \item Slides: \href{aula-05.html}{aula-05.html} e ementa: \href{aula-05.pdf}{aula-05.pdf}.
  \item Livro-base (PDF): \href{../99-Materiais/Engenharia%20de%20software%20projetos%20e%20processos%20by%20Wilson%20de%20P%C3%A1dua%20Paula%20Filho%20.epub.pdf}{Engenharia de software: projetos e processos}.
  \item Vídeos complementares (links nos slides).
\end{itemize}

\section{Cronograma e Descrição Detalhada (120 Minutos)}
\subsection*{Parte 1: Abertura e contexto (10 min)}
\begin{itemize}[leftmargin=*,itemsep=0.35em]
  \item Retomar requisitos (Aula 4) e reforçar que protótipo ajuda a validar entendimento.
  \item Apresentar objetivos e entregáveis da atividade (fluxo + wireframe + perguntas).
\end{itemize}

\subsection*{Parte 2: Conceitos e tipos de protótipo (25 min)}
\begin{itemize}[leftmargin=*,itemsep=0.35em]
  \item Definir prototipagem e fidelidade.
  \item Discutir custo de mudança e por que prototipar cedo é mais barato.
  \item Exemplos: papel vs.\ protótipo navegável (quando escolher).
\end{itemize}

\subsection*{Parte 3: Fluxo de telas e estados (20 min)}
\begin{itemize}[leftmargin=*,itemsep=0.35em]
  \item Montar um fluxo mínimo no quadro (3--5 telas).
  \item Considerar estados: erro, carregando, vazio, sucesso.
  \item Checklist de fluxo: início, decisão, retorno, cancelamento.
\end{itemize}

\subsection*{Parte 4: Wireframe (20 min)}
\begin{itemize}[leftmargin=*,itemsep=0.35em]
  \item Mostrar componentes simples e hierarquia visual.
  \item Orientar: texto curto, botões claros, mensagens úteis.
\end{itemize}

\subsection*{Parte 5: Atividade prática (35 min)}
\begin{itemize}[leftmargin=*,itemsep=0.35em]
  \item Em grupos, escolher 1 funcionalidade do ``app organizador'' (ex.: lembrete de prova).
  \item Entregar:
  \begin{itemize}[leftmargin=*,itemsep=0.25em]
    \item Fluxo mínimo (3--5 telas) desenhado.
    \item Wireframe da tela principal.
    \item 3 perguntas de validação com usuários (o que observar).
  \end{itemize}
  \item Rodada de ``teste de corredor'': um grupo avalia outro com 2 elogios e 1 sugestão.
\end{itemize}

\subsection*{Parte 6: Fechamento (10 min)}
\begin{itemize}[leftmargin=*,itemsep=0.35em]
  \item Consolidar lições: validar cedo reduz risco e retrabalho.
  \item Ponte para a Aula 6: usar requisitos e protótipos para apoiar modelagem (UML).
\end{itemize}

\section{Avaliação}
\begin{itemize}[leftmargin=*,itemsep=0.25em]
  \item Entrega do fluxo, wireframe e perguntas de validação.
  \item Coerência com requisitos e clareza do fluxo.
  \item Participação no feedback cruzado.
\end{itemize}

\section{Referências}
\begin{itemize}[leftmargin=*,itemsep=0.25em]
  \item \href{../99-Materiais/Engenharia%20de%20software%20projetos%20e%20processos%20by%20Wilson%20de%20P%C3%A1dua%20Paula%20Filho%20.epub.pdf}{PAULA FILHO, Wilson de Pádua. Engenharia de software: projetos e processos. 4. ed. 2019.}
\end{itemize}

\end{document}
